\section{Introduction}
\label{sec:introduction}

Joint-embedding predictive architectures and structured world models seek
representations that support prediction and planning without reconstructing
every observed coordinate \citep{assran2023ijepa,assran2024vjepa,
kipf2020contrastive}.  A scalar latent loss, however, does not identify which
part of a structured state failed.  Two predictions with the same latent
error may differ sharply in whether they preserve a typed label, an edge, a
metric comparison, a directed relation, or an order constraint.  This matters
when utility depends on one of those predicates rather than on average latent
fidelity.

This paper studies a narrower question than general world-model superiority:
under a declared finite structural signature, can structural discrepancies
detect errors that survive matched scalar baselines, and do those
discrepancies remain informative under independent world generation,
open-loop composition, and plan selection?  The mathematical interface uses a
single typed correspondence across five layers.  Correspondence-based
distances are established tools for comparing metric and structured objects
\citep{memoli2011gromov,vayer2020fused}; our contribution here is not the
general idea of correspondence.  It is the frozen combination of a coherent
multi-layer state, codebook-free constructive decoding, routing-specific
controls, and failure-preserving sealed evaluation.

The empirical study has four design commitments.
\begin{enumerate}[label=(\roman*)]
    \item Development and sealed worlds are generated from disjoint committed
    seeds, and each sealed gate is executed once.
    \item Six controls share input information, active-parameter count,
    optimization updates, and tuning budget.
    \item Worlds, rather than transitions or tasks, are the independent
    inferential units; optimizer seeds are nested replicates.
    \item Positive results and failed boundaries are reported together.
\end{enumerate}

Three gates are evaluated.  The first tests structural out-of-distribution
(OOD) prediction under unseen mechanisms.  The second tests recursive
open-loop error over horizons up to 32.  The third asks whether task-local
structural discrepancies improve a frozen probabilistic model of plan
failure.  Proper Brier scores provide the probabilistic criterion
\citep{gneiting2007proper}, and a discrete-time hazard represents first-failure
timing \citep{singer1993time}.  Predictor development uses leave-one-world-out
assessment and is frozen before sealed evaluation, following the general
separation between model development and evaluation data
\citep{varma2006bias,collins2024tripod}.

Structural prediction requires a chain from raw observation through entity
discovery, cross-time identity, relations, structural variables, and a
dependency graph to prediction. This study addresses the later stages of that
chain. The entity set and its cross-time identity are supplied by the generator;
what we test is graph recovery and prediction once those are given. Discovering
entities and maintaining their identity from raw observation are outside the
estimand of this paper.

The third gate passes its prespecified primary analysis, but its interpretation
is intentionally limited.  Its discrepancies and regret labels share the same
held-out candidate trajectories; this is a same-task criterion association, not a
deployment-time prognosis that is available without oracle endpoints.
Moreover, a stronger raw layer-aware sensitivity does not reach the
prespecified effect floor, and two temporally separated development
diagnostics fail.  These negative results prevent the scalar-baseline result
from being read as evidence that the full correspondence metric is superior to
every layerwise structural audit.
